\section{Introducción}

En el marco del desarrollo del proyecto de curso de Principios y Tecnologías de Inteligencia Artificial (PTIA), se exploran tres líneas de investigación que presentan retos técnicos significativos y oportunidades claras para la aplicación de modelos avanzados de Inteligencia Artificial (IA). Estas líneas fueron seleccionadas por su capacidad para integrar conocimientos multidisciplinarios, desde el modelado físico hasta la optimización de grafos, y por su relevancia en contextos industriales y académicos reales.

\subsection{Identificación y justificación de problemáticas}

La primera línea de investigación propuesta es la \textbf{predicción inteligente de puntos calientes en un portátil}. El calentamiento excesivo en dispositivos portátiles representa un desafío crítico, ya que reduce el rendimiento operativo y la vida útil de los componentes internos \cite{investigacion_termica}. Aunque la industria emplea materiales con alta conductividad térmica para la disipación, los sistemas pasivos actuales presentan limitaciones físicas ante las cargas de procesamiento modernas \cite{investigacion_termica}. Esta problemática exige un modelado físico-matemático de la transferencia de calor, donde la IA pueda aproximar distribuciones térmicas y recomendar ajustes dinámicos en la frecuencia de la Unidad Central de Procesamiento (CPU) y en los sistemas de ventilación \cite{investigacion_termica}.

En segundo lugar, se plantea el desarrollo de un \textbf{planificador inteligente de ruta de estudio}. El estudiantado universitario suele enfrentarse a trayectorias académicas ineficientes y a la pérdida de motivación debido a la complejidad de las mallas curriculares, que incluyen múltiples prerrequisitos y restricciones temporales \cite{sistema_recomendacion_academica}. La raíz de esta problemática reside en la naturaleza intrincada del grafo académico, donde la secuencia de asignaturas condiciona directamente el tiempo de graduación y el riesgo de retrasos \cite{sistema_recomendacion_academica}. Al modelar el mapa de conocimientos como un \textbf{Grafo Dirigido Acíclico (DAG)}, este proyecto busca identificar rutas que maximicen el dominio temático en el menor tiempo posible, integrando programación lineal e IA para personalizar el flujo académico según el ritmo de aprendizaje real de cada estudiante \cite{sistema_recomendacion_academica}.

Finalmente, se explora la implementación de un \textbf{detector de anomalías en telemetría de Fórmula 1}. Durante una competencia, los equipos generan aproximadamente 1.1 Terabytes (TB) de datos, lo que hace que las anomalías sutiles en componentes críticos como el estrés en el sistema de frenado o defectos en los neumáticos resulten difíciles de detectar mediante supervisión manual \cite{f1_telemetry, fastf1_python}. Mediante el análisis de series de tiempo multivariadas vuelta a vuelta, la IA permite identificar desviaciones del patrón operativo estándar para habilitar mantenimiento predictivo y decisiones estratégicas en tiempo real \cite{f1_predictive_maintenance}. Esta línea resulta especialmente adecuada para el curso, dado el acceso a datos reales públicos y la posibilidad de aplicar técnicas de IA sin supervisión en un entorno altamente dinámico \cite{mercedes_f1_data}.

\pagebreak